\scanpage{33}
\deleted{$(iii)\Longrightarrow(i)$ (cas $R$ réduit)}

\[
(iii)\Longrightarrow\quad (ii')\left\{\begin{gathered}\exists\text{ morphisme }\deleted{\text{un}}\text{ fini surjectif net }S'\longrightarrow S\text{ tel que}\\ \widehat A_{\widehat S'}\text{ soit algébrique.}\end{gathered}\right.
\]

Soit \deleted{$\unsure{N_n=P_n}$}
\[
N_n=\underline{NS}_{A_n/S_n},\qquad P_n=\text{\unsure{sous-schéma} ouvert et fermé défini par les \uline{polarisations}}
\]
\deleted{\unsure{fermé}}

Les $P_n$ se \unsure{recollent} et forment un schéma formel $\widehat P$ sur $\operatorname{Spf}R=\widehat S$, \deleted{somme} \inserted{provenant d'une famille \uuline{dénombrable}} de schémas finis nets $P^\alpha_{\unsure{1}}$ sur $S$
\[
\widehat P=\coprod_\alpha\widehat P^\alpha.
\]
Soit
\[
Z_\alpha=\operatorname{Im}(P^\alpha\longrightarrow S)
\]
L'hyp. (iii) signifie que pour $B$ comme dans l'énoncé, il existe un \inserted{$\alpha$ et un} $S$-morphisme
\[
\deleted{\operatorname{Spec}B'}\quad T'\longrightarrow P^\alpha
\]
donc le pt générique $x$ de $T=\operatorname{Spec}B$ est dans $Z_\alpha$.
Donc
\[
Z=\bigcup_\alpha Z_\alpha
\]
contient tous les $x\in S$ tels que $\dim\overline{\{x\}}=1$.
Par un lemme général, il s'ensuit qu'il y a une suite \uline{finie} des $Z_\alpha$ qui recouvre $S$. Nous prenons pour $S'$ la somme des $P^\alpha$ correspondants.

$\unsure{(iii)|(ii')}\Longrightarrow$ \deleted{$(i)$ (cas $R$ réduit).} \inserted{$(i')$:} $\widehat A_{\widehat S_{\unsure{2|1}}}$ est algébrique, où $S_1=S_{\text{réd}}$.

On peut supposer déjà $S$ réduit. Quitte à faire une extension de la base $S_1\longrightarrow S$ (et redescendre, \deleted{$\unsure{g}$} ce qui est possible par Raynaud), on peut supposer que

\mspage{32}
